\section{Simulación mediante autómatas celulares}

Con el tiempo, se han sugerido diferentes versiones del modelo clásico de autómatas celulares, adaptadas a situaciones sin una estructura clara, que son probabilísticas, diversas o continuas. A continuación, se describen algunas de estas extensiones para el modelado de sistemas complejos.

\begin{itemize}
	
	\item Los autómatas celulares en grafos amplían la idea de la estructura de red clásica al permitir que las vecindades no sean iguales: cada célula puede tener un número diferente de vecinos. Esta flexibilidad los hace adecuados para modelar fenómenos sobre grafos, como redes sociales o flujos en infraestructura. Sin embargo, diseñar manualmente las reglas de transición en estos entornos es una tarea compleja, lo que ha motivado propuestas como los \textit{Graph Neural Cellular Automata}, que permiten aprender dichas reglas automáticamente \cite{neuripsGCA}.
	
	
	
	\item Los autómatas celulares probabilísticos extienden el modelo clásico al incorporar reglas de transición estocásticas. En vez de una evolución que sigue un camino fijo, los estados futuros dependen de probabilidades que se ven afectadas por el entorno. Esto ayuda a entender la incertidumbre, el ruido y los comportamientos que surgen de manera inesperada. Son especialmente útiles para simular dinámicas como propagación de epidemias o sistemas sociales, donde los mecanismos de transición no siempre están completamente definidos \cite{gacsPCA}.
	
	\item Los autómatas celulares no homogéneos permiten que diferentes células tengan diferentes funciones de transición. Para esto, se utilizan dos planos llamados \textit{data plane}, el cual contiene la evolución del sistema; y \textit{control plane}, el cual define las reglas de transición para cada célula. Esto permite estudiar el comportamiento y la interacción entre diferentes reglas que comparten espacio \cite{hartmanIHCA}.
	
	
	
	\item Los autómatas celulares continuos pueden ser continuos en los estados, en el tiempo y en el espacio \cite{nks}:
	
	\begin{itemize}
		
		\item Los estados continuos adquieren valores reales y no solo enteros.
		
		\item El tiempo continuo se define mediante una función que determina el intervalo de espera para actualizar la célula en función de su estado actual.
		
		\item El espacio continuo se obtiene al asignar diferentes pesos a las células dentro de la vecindad \cite{beckerCCA}.
		
	\end{itemize}
	
\end{itemize}

Aunque estas variantes han ampliado el uso de los autómatas celulares, un problema que se repite es definir correctamente las reglas de transición para lograr un fenómeno específico. Este reto ha motivado la introducción de los autómatas celulares neuronales, en los cuales la función de transición no se especifica explícitamente, sino que es aprendida mediante una red neuronal artificial. De este modo, las reglas locales emergen directamente de los datos, lo que facilita modelar dinámicas de alta complejidad.

Un primer avance en esta línea fue presentado por Gilpin \cite{gilpinNCA}, demostrando que una red neuronal convolucional puede aprender el comportamiento de un autómata celular clásico. Su estudio destaca que tanto los autómatas celulares como las redes neuronales convolucionales comparten dos propiedades fundamentales: (1) la actualización basada en una vecindad o \textit{kernel}, y (2) la aplicación simultánea de dicha función en todo el espacio.

En particular, Gilpin entrenó una red neuronal convolucional con un \textit{kernel} de $3 \times 3$ (mismo tamaño que la vecindad de Moore) para imitar el comportamiento del \textit{glider}\footnote{Una de las configuraciones más importantes por su capacidad de moverse en el espacio} en el Juego de la Vida de John Conway. Este trabajo estableció una conexión entre los autómatas celulares y las redes neuronales convolucionales, mostrando que estas redes artificiales pueden actuar como mecanismos de transición local aprendidos.

Posteriormente, esta línea de investigación fue extendida por Mordvintsev et al., quienes exploraron representaciones continuas para morfogénesis artificial (\textit{Growing NCA}), y por Niklasson et al., que aplicaron el enfoque a la generación de texturas dinámicas sobre superficies tridimensionales (\textit{Mesh NCA}). Otros autores también han sugerido usar autómatas celulares neuronales en robótica, clasificación de patrones y simulación de sistemas fotónicos. Esto ha fortalecido su papel como herramienta para modelar sistemas que se autoorganizan.